\section{Notebook-to-OSL method}
\label{sec:method}

Notebook-to-OSL is a workflow for comparing situated and interview-based elicitation, tracing interpretation changes through operator review, and assessing formal transformation. The method is grounded in the alignment in Section~\ref{sec:research-design}: contextual inquiry and situated learning define the situated capture setting for RQ1; a bounded post-session interview provides the comparison source; cognitive task analysis and operator review define the interpretation checks for RQ2; and requirements engineering, MBSE, and OSL define traceability and model-readiness for RQ3. Figure~\ref{fig:notebook-to-osl-workflow} summarises the core transformation path.

\begin{figure}[htbp]
    \centering
    \begin{tikzpicture}[
  node distance=3.5mm,
  box/.style={draw, rounded corners, align=center, minimum width=48mm, minimum height=6mm, font=\scriptsize},
  side/.style={draw, rounded corners, align=center, minimum width=37mm, minimum height=6mm, font=\scriptsize},
  arrow/.style={-{Latex[length=2mm]}, thick}
]

\node[box] (scope) {1. Scope and consent\\training focus, roles, confidentiality};
\node[box, below=of scope] (training) {2. Situated operator training\\contextual/apprenticeship capture};
\node[box, below=of training] (raw) {3. Low-friction digital field notes\\note ID, timestamp, raw text};
\node[box, below=of raw] (select) {4. Select candidate notes\\cue, decision, action, risk, rationale};
\node[box, below=of select] (annotate) {5. Annotate strategy fragments\\schema-based interpretation};
\node[box, below=of annotate] (clarify) {6. CTA-style clarification\\cues, thresholds, exceptions, escalation};
\node[box, below=of clarify] (validate) {7. Operator validation\\accepted, corrected, narrowed, rejected, sensitive};
\node[box, below=of validate] (model) {8. Model-ready candidate\\OSL, EARS, SysML, rule, or knowledge graph};

\draw[arrow] (scope) -- (training);
\draw[arrow] (training) -- (raw);
\draw[arrow] (raw) -- (select);
\draw[arrow] (select) -- (annotate);
\draw[arrow] (annotate) -- (clarify);
\draw[arrow] (clarify) -- (validate);
\draw[arrow] (validate) -- (model);

\node[side, right=7mm of annotate] (trace) {Traceability:\\source, rationale, status,\\uncertainty, responsibility};
\draw[arrow] (trace) -- (annotate);
\draw[arrow] (trace) |- (model);

\end{tikzpicture}

    \caption{Core Notebook-to-OSL transformation path from situated capture to validated model-ready candidate. The RQ1 comparison additionally preserves a separate post-session interview source.}
    \label{fig:notebook-to-osl-workflow}
\end{figure}

The method is organised into nine phases. Table~\ref{tab:workflow} identifies the output and the research question supported by each phase.

\begin{table}[htbp]
\centering
\caption{Notebook-to-OSL workflow phases, outputs, and research-question roles.}
\label{tab:workflow}
\begin{tabularx}{\textwidth}{p{0.25\textwidth}p{0.13\textwidth}X}
\toprule
Phase & RQ & Output \\
\midrule
1. Scope and consent & All & Task focus, stakeholder roles, and confidentiality boundary. \\
2. Situated training & RQ1 & Context-bound operator explanations, demonstrations, and observations. \\
3. Low-friction note capture & RQ1 & Timestamped field notes preserving raw phrasing and task context. \\
4. Post-session interview & RQ1 & Separate interview material covering the same task scope. \\
5. Source comparison and selection & RQ1--RQ2 & Shared and source-exclusive knowledge elements plus candidate fragments. \\
6. Initial annotation & RQ2 & Versioned researcher interpretations linked to their sources. \\
7. CTA-style clarification & RQ2 & Clarified cues, alternatives, thresholds, exceptions, and escalation logic. \\
8. Operator validation & RQ2 & Accepted, corrected, narrowed, rejected, or sensitive fragments. \\
9. OSL transformation assessment & RQ3 & Traceable candidate representations, explicit gaps, and failed mappings. \\
\bottomrule
\end{tabularx}
\end{table}

\subsection{Phase 1: Define scope and consent}

Before data collection, the researcher defines one narrow task scope. The scope may be one machine, operation type, monitoring task, recurring abnormality, quality concern, or decision point. The same scope must be used for situated training and the post-session interview so that RQ1 compares elicitation settings rather than unrelated topics. Confidentiality boundaries should define what may be written, photographed, stored, anonymised, or excluded.

\subsection{Phase 2: Participate in situated training}

During the session, the operator trains the researcher in the work process. The researcher prioritises following the explanation, observing artefacts, and understanding task flow. Clarification questions should remain short and context-sensitive. Questions such as ``what made that matter?'', ``how would you know this is normal?'', or ``what would make you stop?'' are preferable to abstract design questions. The value of the setting is that explanations occur while the operator points, demonstrates, responds, and corrects.

\subsection{Phase 3: Capture low-friction field notes}

The researcher writes short digital field notes during the session. The capture interface should be minimal: note identifier, timestamp, raw text, optional task phase, optional artefact reference, and save action. The researcher should not attempt to complete the full annotation schema during live work. Excessive structuring can disrupt the training and bias what is noticed.

\subsection{Phase 4: Conduct the post-session interview}

After the situated session, the researcher conducts a short conventional interview covering the same task scope. The situated notes should not be shown as prompts during the initial interview pass because doing so would collapse the two sources. The interview should ask the operator to explain relevant cues, decisions, actions, exceptions, risks, and escalation paths from memory. This phase does not assume that either source is superior; it creates the evidence needed to examine how they differ.

\subsection{Phase 5: Compare sources and select candidate fragments}

Situated notes and interview material remain separately identified. The researcher codes both with the same broad categories and records which knowledge elements are shared, situated-only, or interview-only. Candidate fragments are then selected when either source contains or implies an observation, context, interpretation, decision, response, rationale, exception, uncertainty, risk, responsibility, trust concern, or escalation path. This source comparison is the primary analysis for RQ1.

\subsection{Phase 6: Create the initial annotation}

Selected material is annotated using the schema in Section~\ref{sec:schema}. Annotation is an interpretation step, not a mechanical transcription. The initial researcher version must be preserved before clarification or operator review. If the researcher is unsure about the observation, context, interpretation, response threshold, or rationale, the fragment should be marked as uncertain or incomplete rather than silently completed. This version history provides the baseline for RQ2.

\subsection{Phase 7: Apply CTA-style clarification}

For fragments containing hidden reasoning or ambiguous response logic, the researcher applies a bounded CTA-inspired clarification step. Useful prompts include:

\begin{itemize}
    \item What first indicated that this situation was normal or abnormal?
    \item What would a novice likely miss here?
    \item What else could this observation mean?
    \item When would this interpretation be wrong?
    \item Which responses are alternatives rather than a fixed sequence?
    \item What criterion would justify selecting one response over another?
    \item Who should be contacted if the operator is uncertain?
\end{itemize}

Clarification answers remain traceable to the fragment and do not overwrite the initial interpretation.

\subsection{Phase 8: Validate with the operator or domain expert}

The operator reviews the annotated fragment. Validation should test the source interpretation, applicability context, alternatives, response logic, thresholds, risk, exception, responsibility, and confidentiality status. Each fragment is classified as \texttt{accepted}, \texttt{corrected}, \texttt{narrowed}, \texttt{rejected}, or \texttt{sensitive}. The study records both the disposition and the substantive change. These change records answer RQ2; a final accepted label alone is insufficient.

\subsection{Phase 9: Assess OSL transformation}

Only validated or explicitly bounded fragments are assessed for OSL transformation. The purpose is not merely to generate syntactically plausible output. The transformation must preserve source traceability, applicability, observations, interpretations or hypotheses, candidate and selected responses, action roles, expected consequences, uncertainty, evidence, and unresolved gaps where these are present in the source material. A fragment may therefore be classified as coherently represented, represented but explicitly incomplete, partially represented with information loss, or unsupported by the current representation. This assessment answers RQ3.

\subsection{Risk mitigation}

The method has predictable risks. The interview may be influenced by the preceding training session, live notes may be incomplete, and the researcher may over-interpret cues or impose formal categories too early. Operators may also explain differently when observed, and confidential information may appear in either source. Notebook-to-OSL mitigates these risks by preserving source separation, versioning initial interpretations, recording clarification and validation changes, marking uncertainty and sensitivity explicitly, and retaining traceability from each candidate representation back to its empirical source.
